\section{The local tail and a single exceptional neuron}
The scalar-statistic conditioning problem admits a complete description.
Its general probabilistic input is the local large-deviation theorem of
Denisov--Dieker--Shneer. The associated conditional-product phenomenon
was established by Armend\'ariz--Loulakis. We verify the local
large-deviation hypotheses for the critical Gaussian ReLU law and give
a direct conditional-density argument including the overshoot.
This is not a new general
one-big-jump theorem, and the conditioning is not a collection of
individual labels.

Put $\alpha=(d-1)/2>1$ and
\[
 C=\frac{(2\pi)^{-(d+1)/2}}{M(t_*)}
 \frac{\rho^\alpha}{d-1}
 \sum_s\int_{v_s^\perp}e^{-h_s(t)/\rho}\,dt.
 \tag{10}
\]
This constant is positive and finite by the transverse norm property.

\begin{lemma}[Density and marks in the critical tail]\label{lem:localtail}
On $(0,\infty)$ the law of $Z$ under $Q_*$ has a density $f_*$ with
\[
 f_*(z)\sim\alpha C z^{-\alpha-1},\qquad
 Q_*(Z>z)\sim C z^{-\alpha}.
 \tag{11}
\]
Its negative tail is exponentially small. Conditional on $Z=z$ as
$z\to\infty$, the parameter vector approaches one of the $2^r$
critical rays, uniformly among them, and
\[
          \frac{a^2+\|U\|^2}{z}\longrightarrow\frac2\rho
 \quad\text{in probability}.                          \tag{12}
\]
A possible atom at $Z=0$ is harmless and is not excluded.
\end{lemma}
\begin{proof}
Inside a small critical cone write
\[
 a=\sigma_s A,\quad U=Rv_s+u,\quad c=AR,\quad v=A-R,
 \quad t=\sqrt c\,u,
\]
where $A,R>0$, $u\perp v_s$, and $\sigma_sG(v_s)=\rho$.
Let $q=A/\sqrt c=\sqrt{A/R}$. The exact corner identity gives
\[
 Z=\rho c-qh_s(t),\qquad
 D_0:=\tfrac12(a^2+\|U\|^2)-Z/\rho
       =\tfrac12v^2+\frac{\|t\|^2}{2c}+\frac q\rho h_s(t),
\]
with Jacobian $c^{-d/2}/(A+R)$ for $dc\,dv\,dt$.
On this cone $q$ is bounded above and below and
$D_0\ge c_1(v^2+\|t\|)$.
Furthermore
\[
 \frac{\partial q}{\partial c}
       =-\frac{qv}{2c\sqrt{v^2+4c}},\qquad
 \frac{\partial Z}{\partial c}=\rho-h_s(t)\frac{\partial q}{\partial c}.
\]
Shrinking the incoming cone ensures $qh_s(t)\le\eta c$ with
$\eta<\rho$. Thus $\partial Z/\partial c\ge\rho-\eta/2>\rho/2$,
and on $Z=z>0$ one has $z/\rho\le c\le z/(\rho-\eta)$.
The change from $c$ to $z$ is therefore regular.

For fixed $(v,t)$ and $z\to\infty$, one has
$c\sim z/\rho$, $q\to1$, $(A+R)/\sqrt c\to2$, and
$\partial Z/\partial c\to\rho$. After multiplying the resulting
density by $z^{\alpha+1}$ its integrand converges to
\[
 \frac{(2\pi)^{-(d+2)/2}}{M(t_*)}
       \frac{\rho^\alpha}{2}
                      e^{-v^2/2-h_s(t)/\rho}.
\]
It is dominated on the whole cone by a constant times
$e^{-c_1(v^2+\|t\|)}$. The cone domains exhaust the transverse
space and the real $v$ axis. Dominated convergence, followed by the
Gaussian integral in $v$, gives exactly $\alpha C$ in (11).

For clarity, the part outside these finitely many cones is negligible
also at the density level. The homogeneous quadratic exponent $D_0$
has a positive gap there on the unit sphere. Integrate the readout at
$a=z/G(U)$; the density integrand is $|G(U)|^{-1}e^{-D_0}$ times
its normalizing constant. The gap bounds it by
\[
 e^{-c_2z}\,|G(U)|^{-1}
           e^{-c_3\{\|U\|^2+z^2/G(U)^2\}}.
\]
Here $a^2+\|U\|^2\ge2z/\rho$. The inequality
$y^{-1}e^{-c_3z^2/y^2}\le C_3/z$ for $z>0$, $y>0$ bounds the
remaining integral by $C_4/z$. Points with $G(U)=0$ contribute only
to the atom at zero. This proves the asserted density asymptotic.
Integration gives the tail. Readout symmetry before tilting implies
$f_*(-z)=e^{-2t_*z}f_*(z)$ on $z>0$, proving the negative-tail claim.
The same dominated calculation makes $(v,t)$ tight given large $z$,
which yields (12) and convergence to the critical rays. Relevant input
sign changes, compensated in the readout, preserve $Z$ and act
transitively on those rays; their limiting probabilities are equal.
\end{proof}

For example, at $d=r=4$, orthonormal transverse coordinates give
\[
 h_s(t)/\rho=\sum_{i=1}^3\max(|b|/2,|t_i|),\qquad
 \int e^{-h_s(t)/\rho}dt
       =32\int_0^\infty(1+u)^3e^{-3u}du=832/27.
\]
The earlier importance estimates $M(8)\approx1.178$ and
$m_c\approx0.1805$ give $C\approx0.0623$; these displayed decimal
values are estimates, not interval certificates.

\begin{theorem}[Scalar conditioning above the critical mean]\label{thm:scalar}
Fix $d\ge4$, an even $2\le r\le d$, and $s>m_c$. Let
$\delta=s-m_c$ and $S_N=\sum_{i=1}^N Z_i$ for independent prior
neurons. Then
\[
 P_0^N(S_N\ge Ns)
 \sim e^{-NI(s)}\frac{\alpha C}{t_*\delta^{\alpha+1}}N^{-\alpha}.
 \tag{13}
\]
Under this conditional law, the overshoot $S_N-Ns$ converges in
law to an exponential variable of rate $t_*$. After moving the neuron
with largest $Z_i$ to the last position, the joint law of all other
$N-1$ parameter vectors converges in total variation to $Q_*^{N-1}$.
For the largest neuron $J_N$,
\[
 Z_{J_N}/N\to\delta,\qquad
 (a_{J_N}^2+\|U_{J_N}\|^2)/N\to2\delta/\rho,
 \qquad \max_{i\ne J_N}\|\theta_i\|/\sqrt N\to0
 \tag{14}
\]
in conditional probability. Its ray is asymptotically uniform among
the $2^r$ critical rays. In particular exactly one neuron exceeds
$\epsilon\sqrt N$ in norm with probability tending to one whenever
$0<\epsilon<\sqrt{2\delta/\rho}$.
\end{theorem}
\begin{proof}
Apply the cited local large-deviation result to the centered critical variable
$X=Z-m_c$. Lemma~\ref{lem:localtail} gives a regularly varying positive
local density and an exponentially bounded negative tail. Choose
$1<\kappa<\min\{\alpha,2\}$. Then $\E_{Q_*}|X|^\kappa<\infty$;
for each fixed interval length $h>0$ its local probability is almost
decreasing, regularly varying, and insensitive to shifts of order
$x^{1/\kappa}=o(x)$. These verify Corollary 2.1 of
\href{https://arxiv.org/abs/math/0703265}{Denisov--Dieker--Shneer (2008)}.
It supplies, uniformly for $y\ge N\delta$,
\[
 Q_*^N\{S_N-Nm_c\in(y,y+h]\}
       \sim Nh f_*(y).
 \tag{15}
\]
The constant mean shift in the one-step density does not alter the
asymptotic. The product-bulk conclusion is consistent with Theorem 1 of
\href{https://arxiv.org/abs/0912.1516}{Armend\'ariz--Loulakis (2011)};
the direct density argument below additionally tracks the overshoot.

Here are details of the further exponential reweighting, which is
necessary because the desired conditional law is under $P_0$, not $Q_*$.
Changing measure gives the residual factor
\[
 D_N=\E_{Q_*^N}\left[e^{-t_*(S_N-Ns)}\mathbf1_{\{S_N\ge Ns\}}\right],
 \qquad P_0^N(S_N\ge Ns)=e^{-NI(s)}D_N.
\]
Partition the nonnegative overshoot axis into intervals of length $h$.
For each bounded overshoot (15) gives probability asymptotic to
$Nh f_*(N\delta)$. The same uniform estimate and regular variation
bound every interval probability by $C_hNh f_*(N\delta)$ for all
large $N$. Thus exponential weights give a summable dominating sequence.
Upper and lower endpoint Riemann sums, with $N\to\infty$ first and
$h\downarrow0$ afterward, yield
\[
 D_N\sim Nf_*(N\delta)\int_0^\infty e^{-t_*u}\,du
                 =Nf_*(N\delta)/t_*.
\]
This proves (13). Starting the same sums at a fixed overshoot $b\ge0$
gives limiting survival probability $e^{-t_*b}$.

Here is the direct product argument. Let $\mu$ be the law of $Z$
under $Q_*$. Given a proposed bulk $z=(z_1,\ldots,z_{N-1})$, put
$b(z)=Ns-\sum_i z_i$ and $m(z)=\max_i z_i$.
After swapping the maximum to the last position, the joint conditional
law of this scalar bulk and the overshoot $u\ge0$ has density
\[
 \frac{N}{D_N}e^{-t_*u}f_*(b(z)+u)
                  \mathbf1_{\{b(z)+u>m(z)\}}             \tag{15a}
\]
with respect to $\mu^{N-1}(dz)\,du$. The maximum is positive and
unique almost surely on the conditioning event, so exchangeability
gives exactly the factor $N$, even if $\mu$ has an atom at zero.
Under the reference probability law
$\mu^{N-1}\otimes\operatorname{Exp}(t_*)$, one has
$b(z)/N\to\delta$, $m(z)/N\to0$, and $u/N\to0$ in probability.
The finite first moment gives the first two claims. Regular variation
of $f_*$ and $D_N\sim Nf_*(N\delta)/t_*$ show that the density ratio
in (15a) relative to this reference law tends to one in probability.
Its integral is one. For any such nonnegative density ratio $R_N$,
$\E|R_N-1|=2\E(1-R_N)_+\to0$ by bounded convergence in probability.
This proves the stronger joint total-variation convergence of bulk and
overshoot to the indicated independent product.
The conditional parameter law given each $Z_i$ is identical under
$P_0$ and $Q_*$; applying these independent mark kernels cannot increase
total variation, so the same conclusion holds for the parameter vectors.

The finite first moment gives the bulk law of large numbers. Combined
with tight overshoot, it implies $Z_{J_N}/N\to s-m_c$.
Lemma~\ref{lem:localtail} gives the marked large-neuron limits in (14).
The critical parameter second moment is finite for $d>3$, by the same
ray-coordinate integral as the first moment of $Z$. Thus
$NQ_*(\|\theta\|>\epsilon\sqrt N)\to0$ for each fixed $\epsilon>0$.
A union bound under the product bulk law, and its total-variation
approximation, prove the final limit in (14).
\end{proof}

The theorem completes the scalar-conditioning application left open in
the first-round draft. Individual-label conditioning is different: full
cube margins require compatible contributions in all antipodal directions,
which is why the later full-output companions have exponentially many
exceptional neurons rather than one. None of these scalar-conditioning
statements asserts an OOD mistake lower bound.

\section{A large parity statistic need not classify the parity}
The scalar-conditioned network has an explicit limiting output, which
also makes the limitation of this conditioning visible. Write
$\phi_v(x)=\relu(v\cdot\bar x)$ and
$\sigma_s=\operatorname{sgn}G(v_s)$. For the original normalization
$F_N=N^{-1}\sum_i\Phi(\theta_i)$, the preceding theorem gives
\[
 F_N\ \Longrightarrow\ m_c\chi+
             \frac{\delta}{\rho}\sigma_S\phi_{v_S}
 \quad\text{in }\mathbb R^{\{-1,1\}^d},                 \tag{16}
\]
where $S$ is uniform among the $2^r$ critical orientations. The first
term follows from the finite-mean product bulk law; homogeneity and
(12) give the second. Thus one posterior draw has a random limiting
function even though the ensemble is invariant under all source-label
symmetries.

\begin{corollary}[Nonvanishing classification error of a scalar-conditioned draw]
Let $r\ge4$ be even and $d\ge r$. If
\[
                 \delta>\frac{m_c\rho\sqrt r}{4},       \tag{17}
\]
the uniform-input classification error of one scalar-conditioned
posterior draw converges in probability to
\[
 q_r=\frac14-\frac1{2^{r+1}}\binom r{r/2}.
 \tag{18}
\]
In particular $q_r\to1/4$ as $r\to\infty$, after the width limit.
For every fixed input, the posterior probability that the draw
misclassifies it also converges to $q_r$. Consequently posterior
majority classification is eventually correct at every input, while a
sampled network retains error (18). Its posterior mean output converges
to $s\chi$.
\end{corollary}
\begin{proof}
Put $r=2m$ and $y_i=S_ix_i$ for $i\le r$. The signed margin in (16)
is
\[
 m_c+\frac\delta\rho(-1)^{m-1}\chi(y)
             \relu\left(\frac{\sum_i y_i}{\sqrt r}\right).
\]
Here $m_c>0$: the untilted statistic is symmetric and nondegenerate,
so its finite positive-endpoint tilted mean is strictly positive.
If $k$ coordinates of $y$ are negative, the second term can be negative
only when $k<m$ and $k\equiv m\pmod2$. In that case $m-k\ge2$,
so its absolute value is at least $4\delta/(\rho\sqrt r)>m_c$.
At every other input the margin is at least $m_c$. All limiting margins
are therefore bounded away from zero. Exactly
\[
 \sum_{\substack{k<m\\k\equiv m\ (2)}}\binom r k
       =2^{r-2}-\frac12\binom r m
\]
inputs in the relevant cube have a negative margin; the identity follows
by pairing $k$ with $r-k$ in the parity-restricted binomial sum.
This count is independent of $S$, so (16) proves convergence of the
sampled classification error. For fixed $x$, uniform $S$ makes $y$
uniform, proving the posterior-probability assertion and, since
$q_r<1/2$, the majority assertion. The finite cube permits the latter
conclusion simultaneously at all inputs for sufficiently large widths.

Finally the exact conditional symmetries force
$\E[F_N\mid S_N\ge Ns]=c_N\chi$, with
$c_N=\E[S_N/N\mid S_N\ge Ns]$.
The uniform interval bound in the proof of (13), summed with exponential
weights, bounds the conditional overshoot tail by $C e^{-t_*b}$
uniformly for all large $N$ and $b\ge0$. Its means are bounded, hence
$c_N\to s$.
\end{proof}

This is a failure of learning from only a global parity statistic, not
a failure of the ordinary posterior conditioned on labeled examples.
It also separates sampling a network from taking the posterior majority;
the latter has no classification failure in this example.
